% page 636 du fac-simile (image p-635 du scan)
% bloc 172.7 x 263.6 mm ; marge gauche 8.0 mm
\pgc{-12.700mm}{0.000mm}{
\l{}
\l{}
\l{}
\l{}
\l{}
\l{}
\l{}
\l{}
\l{\cel{11}628}
\l{}
\l{\cel{7}zinkografio. Procedo analoga a litografado, ma en qua la}
\l{\cel{10}petroplako litografala remplasesas da zinkoplako. -}
\l{}
\l{\cel{7}\sou{zirkono.(kemio}).Korpo simpla, metala, kelke simila a tita-}
\l{\cel{10}no. Atom-pezo : 90,6. -\cel{1}\sou{Simbolo kemiala}\cel{1}: Zr.}
\l{}
\l{\cel{7}\sou{zodiako}. (\sou{astron.}) Zono cielala qua kovras 8 o 9 gradi ye}
\l{\cel{10}singla latero di ekliptiko, ed embracas omna situeso di}
\l{\cel{10}la planeti til nun deskovrita - e dividita aden 12 parti}
\l{\cel{10}qui nomizesis segun la stelaro maxim vicina a li, e quin}
\l{\cel{10}onu nomas la 12 signi di zodiako. - DEFIRS.}
\l{}
\l{\cel{7}\sou{zono}. I. (\sou{geom.}) Parto de la areo di sfero qua formacas qua-}
\l{\cel{10}za bendo kurva inter du plani interparalela. - II. (\sou{geog}.)}
\l{\cel{10}Singla de la granda regioni di la Terglobo, quin separas}
\l{\cel{10}cirkli imaginala, paralela ye equatoro, e di qui la kli-}
\l{\cel{10}mato dependas de lia situeso kompare a Suno. - III. Bendo}
\l{\cel{10}de stofo, de ledro, e c., destinata a klemar la vesti}
\l{\cel{10}an la tayo. - IV. To quo cirkondas quale vesto-zono. -}
\l{\cel{10}DEFIRS.}
\l{}
\l{\cel{7}\sou{zoocecidio}.\cel{2}Deformuro genitata da animalo o planto vivanta.}
\l{}
\l{\cel{7}\sou{zoofito}. Nomo di la animali inferiora, vicina ye la planti,}
\l{\cel{10}qui formacas la brancho quaresma di la animal-regno (de}
\l{\cel{10}qua onu deprenis nia-epoke la protozoi por formacar kin-}
\l{\cel{10}esma). - DEFIRS.}
\l{}
\l{\cel{7}\sou{zoografio}. Deskripto di la animali. - EFIRS.}
\l{}
\l{\cel{7}\sou{zoologio}. Parto di la historio naturala qua traktas la ani-}
\l{\cel{10}mali. - DEFIRS.}
\l{}
\l{\cel{7}\sou{zoospermo}. (sinonimo di \textquotedbl{}spermatozoido\textquotedbl{}. - DEFIS.}
\l{}
\l{\cel{7}\sou{zoosporo}. (\sou{zool.})\cel{2}Sporo (korpuskulo genitala) quan karak-}
\l{\cel{10}terizas cilii vibranta, che ula algi. - DEFIS.}
\l{}
\l{\cel{7}\sou{zuavo}. (\sou{armeo}). En Francia, soldato ek korpo de infantrio}
\l{\cel{10}lejera, uzata en la armeo di Afrika. - DEFIRS.}
\l{}
\l{\cel{7}\sou{zumar}. (\sou{netrans}.) I. Genitar bruiso kontinua, grava e mati-}
\l{\cel{10}da, aludante insekto, o la murmuro konfuza ek la voci}
\l{\cel{10}di plura personi asemblita. - II. Resonigar klosho per}
\l{\cel{10}igar la frapilo shokar alterne dextre e sinistre, ma sen}
\l{\cel{10}ociligar ta klosho. - III. (\sou{medic}.)Audar en sua oreli}
\l{\cel{10}la sono quan genitas ibe la palpitado tro forta da la}
\l{\cel{10}arterii, la stretesko di la tubo audala.}
\l{}
\l{}
\l{\cel{31}---------}
}
